\documentclass[12pt, a4paper]{article}

% ---- 字体与编码 ------------------------------------------------------------
\usepackage{fontspec}
\usepackage{xeCJK}
\setCJKmainfont[BoldFont=SimHei]{SimSun}
\setCJKsansfont{SimHei}
\setCJKmonofont{FangSong}
\setmainfont{Times New Roman}
\setmonofont[
    Path = C:/Windows/Fonts/,
    Extension = .ttf,
    UprightFont = *,
    BoldFont = *b,
    ItalicFont = *i,
    BoldItalicFont = *z
]{consola}
\newfontfamily\codefont[
    Path = C:/Windows/Fonts/,
    Extension = .ttf,
    UprightFont = *,
    BoldFont = *b,
    ItalicFont = *i,
    BoldItalicFont = *z
]{consola}

% ---- 页面布局 --------------------------------------------------------------
\usepackage[top=2.5cm, bottom=2.5cm, left=3cm, right=3cm]{geometry}
\usepackage{setspace}
\onehalfspacing
\emergencystretch=2em

% ---- 数学、图表与引用 ------------------------------------------------------
\usepackage{amsmath, amssymb}
\usepackage{booktabs}
\usepackage{graphicx}
\usepackage{float}
\usepackage{caption}
\captionsetup{font=small, labelfont=bf}
\renewcommand{\contentsname}{目录}
\renewcommand{\tablename}{表}
\renewcommand{\figurename}{图}

% ---- 代码块 ----------------------------------------------------------------
\usepackage{listings}
\usepackage{xcolor}
\definecolor{codebg}{RGB}{246,248,250}
\definecolor{keywordcolor}{RGB}{207,34,46}
\definecolor{commentcolor}{RGB}{106,115,125}
\definecolor{stringcolor}{RGB}{10,132,255}
\definecolor{numbercolor}{RGB}{140,149,159}
\definecolor{accentbar}{RGB}{9,105,218}
\renewcommand{\lstlistingname}{代码片段}
\lstset{
    backgroundcolor=\color{codebg},
    basicstyle=\codefont\small\linespread{1.1}\selectfont,
    keywordstyle=\color{keywordcolor}\bfseries,
    commentstyle=\color{commentcolor}\itshape,
    stringstyle=\color{stringcolor},
    breaklines=true,
    breakatwhitespace=false,
    frame=l,
    framerule=3pt,
    rulecolor=\color{accentbar},
    framesep=8pt,
    xleftmargin=14pt,
    numbers=left,
    numberstyle=\tiny\color{numbercolor},
    numbersep=8pt,
    tabsize=4,
    columns=fullflexible,
    keepspaces=true,
    showstringspaces=false,
    aboveskip=10pt,
    belowskip=6pt,
}

% ---- 超链接 ----------------------------------------------------------------
\usepackage{hyperref}
\hypersetup{colorlinks=true, linkcolor=blue, urlcolor=blue}
\usepackage{cleveref}
\crefformat{table}{#2表~#1#3}
\crefformat{figure}{#2图~#1#3}
\crefformat{equation}{#2式~(#1)#3}

\title{基于 E203 RISC-V 软硬件协同的 INT8 GEMV 加速器设计}
\author{陈炯涛 董山}
\date{2026 年 5 月}

\renewcommand{\theequation}{\thesection.\arabic{equation}}
\makeatletter
\@addtoreset{equation}{section}
\makeatother

\begin{document}
\maketitle
\tableofcontents
\newpage
\setcounter{section}{-1}

\section{快速预览简介}
\label{sec:overview}

\paragraph{项目目标}
本项目面向大模型推理中的线性层计算，基于 HummingBird E203 RISC-V 处理器和 ALINX AXU3EG FPGA 开发板，实现一个 8$\times$8 INT8 GEMV 矩阵向量乘加硬件加速器。项目采用先软件基线、再硬件外设、最后上板 ILA 与 PL UART 验证的流程，形成从 E203 仿真到 FPGA 板端运行的完整软硬件协同验证链路。

\paragraph{核心技术}
\begin{itemize}
    \item HummingBird E203 RISC-V SoC，使用 ICB peripheral bus 接入自定义 memory-mapped 外设
    \item 8$\times$8 INT8 GEMV 软件基线，基于 E203 \texttt{mcycle} CSR 统计 kernel cycle
    \item \texttt{gemv\_accel} 硬件加速器，采用 8 路并行 signed MAC，每周期处理一列
    \item Vivado 2022.2 FPGA 工程迁移，目标器件为 \texttt{xczu3eg-sfvc784-1-i}
    \item ILA 在线调试，重点观察 GEMV 外设内部状态、计算列计数和完成信号
    \item PL UART 串口输出，获取完整 \texttt{y[0..7]} 正确性结果和端到端 cycle 数据
\end{itemize}

\paragraph{当前重要指标}
\begin{itemize}
    \item 板上软件基线：8$\times$8 INT8 GEMV kernel 在 E203 上耗时 1630 cycles
    \item 板上硬件端到端：CPU 通过 MMIO 调用 GEMV 加速器耗时 243 cycles
    \item 端到端加速比：$1630 / 243 \approx 6.71\times$
    \item RTL 推导硬件核心延迟：8 个 COMPUTE 周期 + 1 个 WRDONE 周期，共 9 cycles
    \item 外设基地址：\texttt{0x1004\_1000}
    \item 上板结果：PL UART 输出显示 \texttt{SW GEMV PASS} 和 \texttt{HW GEMV PASS}
\end{itemize}

\paragraph{当前状态}
软件仿真链路已经完成，硬件加速器已经集成到 E203 SoC 并完成 FPGA 上板验证。BRAM 推断失败导致 CPU 卡死在 PC 0 的关键问题已经定位并修复。最终板端测试中，E203 通过 PL UART 输出完整软件和硬件 GEMV 结果，8 个输出元素均与 golden values 一致，并获得软件 baseline 与硬件端到端 cycle 数据。

\section{项目背景与总体目标}

\subsection{背景}

大语言模型推理中存在大量线性层计算，例如 Q/K/V projection、attention output projection 和 FFN 层。在线性层中，矩阵向量乘法（GEMV）和矩阵矩阵乘法（GEMM）占据主要计算量。为了降低存储带宽、访存能耗和乘加资源开销，实际推理系统通常使用 INT8 或更低比特量化数据。

本课程设计选择 INT8 GEMV 作为硬件加速对象。一方面，GEMV 算子形式简单，便于在 E203 软核系统中建立完整的正确性验证；另一方面，GEMV 具有清晰的数据复用和并行乘加结构，适合作为 FPGA 数据通路设计的入门加速模块。

\subsection{目标}

项目目标不是单独完成一个 RTL 模块，而是完成一条可运行的软硬件协同链路：

\begin{enumerate}
    \item 在 WSL 中搭建 E203 iVerilog 软件仿真环境，确认用户 C 程序可以在 E203 RTL 模型上运行。
    \item 实现 8$\times$8 INT8 GEMV 软件基线，验证输出正确性并测量 E203 cycle 数。
    \item 在 Vivado 工程中实现 GEMV hardware accelerator，并将其接入 E203 peripheral ICB bus。
    \item 编写固件，通过 MMIO 配置加速器、启动计算、轮询完成状态并读取输出。
    \item 在 FPGA 板端通过 ILA 和 PL UART 观察 CPU、加速器与程序输出，验证完整软硬件链路。
\end{enumerate}

\section{开发环境与工程结构}

\subsection{硬件与软件环境}

本项目使用的目标开发板为 ALINX AXU3EG，FPGA 型号为 \texttt{xczu3eg-sfvc784-1-i}。主要工具和版本如\cref{tab:env} 所示。

\begin{table}[H]
    \centering
    \caption{项目开发环境}
    \label{tab:env}
    \begin{tabular}{ll}
        \toprule
        \textbf{工具或平台} & \textbf{版本或说明} \\
        \midrule
        FPGA 开发板 & ALINX AXU3EG \\
        FPGA 器件 & \texttt{xczu3eg-sfvc784-1-i} \\
        Vivado & 2022.2 \\
        RISC-V GCC & Nuclei GCC 9.2.0 \\
        RTL 仿真 & Icarus Verilog 12.0 \\
        SDK & hbird-sdk \\
        板端调试 & Vivado Hardware Manager + ILA，PL UART 115200 8N1 \\
        \bottomrule
    \end{tabular}
\end{table}

\subsection{工程结构}

课程提供的主要工程目录位于 \texttt{lecture2\_0420\_release/}。其中，Vivado 工程、固件工程和辅助工具目录如\cref{tab:layout} 所示。

\begin{table}[H]
    \centering
    \caption{主要工程目录}
    \label{tab:layout}
    \begin{tabular}{ll}
        \toprule
        \textbf{路径} & \textbf{作用} \\
        \midrule
        \texttt{lecture2\_0420\_release/trans/} & Vivado 2022.2 FPGA 工程 \\
        \texttt{trans.srcs/sources\_1/new/} & 新增 RTL 模块目录 \\
        \texttt{test1/application/main.c} & FPGA 固件主程序 \\
        \texttt{hexdump-2.1.0/} & hex/bin 转换和辅助脚本目录 \\
        \texttt{/home/cjt/csapp/e203\_hbirdv2/} & WSL 中的 E203 iVerilog 仿真工程 \\
        \texttt{/home/cjt/csapp/hbird-sdk/} & WSL 中的 SDK 应用工程 \\
        \bottomrule
    \end{tabular}
\end{table}

\section{软件基线设计与仿真结果}

\subsection{软件 GEMV kernel}

软件基线程序实现固定尺寸 8$\times$8 INT8 GEMV：

\begin{equation}
    y_i = \sum_{j=0}^{7} W_{i,j} x_j,\quad i=0,1,\dots,7
\end{equation}

其中矩阵 $W$ 和输入向量 $x$ 使用 \texttt{int8\_t}，累加器和输出 $y$ 使用 \texttt{int32\_t}。软件 kernel 采用双重循环实现，每个输出元素需要 8 次 signed 乘加，总共完成 64 次乘加操作。

该阶段的重点是建立可复现的软件基线：一方面确认 INT8 GEMV 的 C 程序能够在 E203 RTL 仿真环境中正确运行，另一方面测量纯软件实现的 cycle 数，为后续硬件加速效果提供参考。

\subsection{软件基线在整体框架中的作用}

软件基线不是最终优化对象，而是整个设计框架中的参考坐标。它承担三个作用：第一，提供一组已经验证过的输入、输出和 golden values，保证后续硬件结果可以逐项比较；第二，给出 E203 顺序执行 INT8 GEMV 时的 cycle 数，作为加速比计算的分母；第三，保留完全由 C 程序实现的路径，便于区分“算法本身错误”和“硬件外设或总线接口错误”。

软件程序在 hbird-sdk 的 baremetal 环境中编译为 E203 可执行镜像，再加载到 E203 RTL 仿真模型中执行。报告中不展开具体命令，因为这些命令只属于工具链复现步骤；从设计角度看，关键是软件基线与硬件加速器使用同一类数据格式：矩阵和输入向量为 \texttt{int8\_t}，输出为 \texttt{int32\_t}，从而保证软硬件比较的含义一致。

仿真程序会打印每个输出元素，并将结果与 golden values 比较。早期软件基线结果如\cref{tab:sw-result} 所示。

\begin{table}[H]
    \centering
    \caption{8$\times$8 INT8 GEMV 软件基线结果}
    \label{tab:sw-result}
    \begin{tabular}{ccc}
        \toprule
        \textbf{输出元素} & \textbf{软件结果} & \textbf{Golden} \\
        \midrule
        $y_0$ & 56 & 56 \\
        $y_1$ & -14 & -14 \\
        $y_2$ & 25 & 25 \\
        $y_3$ & -23 & -23 \\
        $y_4$ & 33 & 33 \\
        $y_5$ & 24 & 24 \\
        $y_6$ & 36 & 36 \\
        $y_7$ & -8 & -8 \\
        \bottomrule
    \end{tabular}
\end{table}

E203 iVerilog 仿真输出显示 \texttt{INT8 GEMV PASS}，并测得早期软件基线：

\begin{equation}
    T_{\text{software}} = 1638 \text{ cycles}
\end{equation}

该结果作为当前项目的软件 baseline。

后续 FPGA 板端测试使用对角矩阵数据集，PL UART 实测软件 GEMV 为 1630 cycles。两个数值来自不同测试数据和运行环境，最终性能分析以板端 UART 输出的 1630 cycles 作为软件端到端对比基准。

\section{硬件加速器设计}

\subsection{总体方案}

硬件加速器被设计为 E203 SoC 的 memory-mapped ICB 外设。CPU 通过普通 load/store 指令访问外设寄存器，不需要修改 E203 core 本身。整体运行流程如下：

\begin{enumerate}
    \item CPU 将 8$\times$8 INT8 矩阵 $W$ 写入外设内部 \texttt{W\_mem}。
    \item CPU 将 8 维 INT8 向量 $x$ 写入外设内部 \texttt{x\_mem}。
    \item CPU 向 \texttt{CTRL} 寄存器写 start bit。
    \item 加速器在硬件中执行 8 路并行 signed MAC。
    \item 加速器置位 done bit，CPU 轮询后读取 \texttt{Y\_MEM}。
\end{enumerate}

该方式的优点是接口简单、调试直观，并且与 E203 原有 peripheral bus 结构兼容。

\subsection{寄存器映射}

加速器 RTL 文件为：

\begin{lstlisting}[caption={GEMV 加速器 RTL 文件路径}]
lecture2_0420_release/trans/trans.srcs/sources_1/new/gemv_accel.v
\end{lstlisting}

外设基地址为 \texttt{0x1004\_1000}，寄存器映射如\cref{tab:regmap} 所示。

\begin{table}[H]
    \centering
    \caption{GEMV 加速器 MMIO 寄存器映射}
    \label{tab:regmap}
    \begin{tabular}{p{3.0cm}p{4.0cm}p{6.6cm}}
        \toprule
        \textbf{偏移地址} & \textbf{名称} & \textbf{说明} \\
        \midrule
        \texttt{0x000} & \texttt{CTRL} & bit[0] 为 start，bit[1] 为 done \\
        \texttt{0x010--0x04C} & \texttt{W\_MEM} & 8$\times$8 INT8 权重矩阵，row-major，4 个 int8 打包为 1 个 word \\
        \texttt{0x050--0x054} & \texttt{X\_MEM} & 8 维 INT8 输入向量，4 个 int8 打包为 1 个 word \\
        \texttt{0x060--0x07C} & \texttt{Y\_MEM} & 8 个 int32 输出元素，只读 \\
        \texttt{0x080} & \texttt{DEBUG\_SW\_CYCLES} & 固件写入的软件 GEMV cycle delta \\
        \texttt{0x084} & \texttt{DEBUG\_HW\_CYCLES} & 固件写入的硬件辅助端到端 cycle delta \\
        \bottomrule
    \end{tabular}
\end{table}

\subsection{计算数据通路}

加速器内部包含 8 个 signed 32-bit accumulator，分别对应 8 个输出元素。每个计算周期处理一列输入：当前列索引为 \texttt{col}，硬件读取输入向量的第 \texttt{col} 项和权重矩阵第 \texttt{col} 列，并并行执行：

\begin{equation}
    acc_i \leftarrow acc_i + W_{i,col} \times x_{col},\quad i=0,1,\dots,7
\end{equation}

FSM 状态如\cref{tab:fsm} 所示。

\begin{table}[H]
    \centering
    \caption{GEMV 加速器 FSM}
    \label{tab:fsm}
    \begin{tabular}{ll}
        \toprule
        \textbf{状态} & \textbf{作用} \\
        \midrule
        \texttt{IDLE} & 等待 CPU 写 start bit \\
        \texttt{COMPUTE} & 连续处理 8 列输入，每列执行 8 路并行 MAC \\
        \texttt{WRDONE} & 将 accumulator 写入 \texttt{y\_mem[0..7]} 并置位 done \\
        \bottomrule
    \end{tabular}
\end{table}

因此硬件核心计算延迟为：

\begin{equation}
    T_{\text{core}} = 8 + 1 = 9 \text{ cycles}
\end{equation}

与软件基线相比，核心计算理论加速比为：

\begin{equation}
    S_{\text{core}} = \frac{1630}{9} \approx 181
\end{equation}

需要强调的是，$T_{\text{core}}$ 只统计硬件内部乘加和写回结果的时间，不包含 CPU 写入矩阵、写入向量、启动外设、轮询 done 和读取结果的 MMIO 开销。因此最终展示时应同时给出硬件核心延迟和硬件辅助端到端延迟。

\section{SoC 集成与固件设计}

\subsection{E203 外设总线集成}

本设计没有修改 E203 core 内部流水线，而是把 GEMV 加速器作为一个 memory-mapped peripheral 接入 SoC。这样做的原因是接口边界清晰：CPU 仍然执行标准 load/store 指令，硬件加速器只需要遵守 E203 peripheral ICB bus 的握手协议即可。相比在执行级直接加入新指令，这种方式对原处理器侵入更小，也更适合课程工程中的 FPGA 验证。

原工程中 O14 slot 对应 \texttt{0x1004\_1000--0x1004\_1FFF} 地址窗口，原用途是 AXI example peripheral。本项目保留该地址窗口和总线路由关系，将其后端替换为直接 ICB slave 形式的 \texttt{gemv\_accel}。因此 CPU 访问 \texttt{0x1004\_1000} 起始的寄存器时，总线 fabric 会把请求送到 GEMV 外设，而不是送往原示例 AXI 外设。

这个集成方式形成了三层结构：E203 core 负责运行控制程序，ICB bus 负责地址译码和读写握手，\texttt{gemv\_accel} 负责保存输入数据、执行乘加和返回输出。三层之间只通过寄存器读写交互，避免了软件和硬件实现细节互相耦合。

\subsection{固件 MMIO 访问}

固件是软硬件协同框架中的控制层。它不参与硬件内部 MAC 的实现，而是负责把软件中的矩阵和向量打包为 32-bit word，通过 MMIO 写入外设寄存器，然后写 \texttt{CTRL.start} 启动计算，轮询 \texttt{CTRL.done}，最后读取 8 个 32-bit 输出结果。

固件中定义了与硬件一致的寄存器地址。该定义相当于软硬件接口契约：只要寄存器偏移、数据打包顺序和状态位语义保持一致，CPU 侧代码和 RTL 内部实现就可以独立演进。

\begin{lstlisting}[language=C,caption={固件中的 GEMV MMIO 定义}]
#define GEMV_BASE       0x10041000UL
#define GEMV_CTRL       (*(volatile uint32_t *)(GEMV_BASE + 0x000))
#define GEMV_W_REG(k)   (*(volatile uint32_t *)(GEMV_BASE + 0x010 + (k)*4))
#define GEMV_X_REG(k)   (*(volatile uint32_t *)(GEMV_BASE + 0x050 + (k)*4))
#define GEMV_Y_REG(i)   (*(volatile uint32_t *)(GEMV_BASE + 0x060 + (i)*4))
#define GEMV_DBG_SW_CYCLES (*(volatile uint32_t *)(GEMV_BASE + 0x080))
#define GEMV_DBG_HW_CYCLES (*(volatile uint32_t *)(GEMV_BASE + 0x084))
\end{lstlisting}

当前板端验证使用对角矩阵测试数据：

\begin{equation}
    W_{i,i}=i+1,\quad x=\{1,2,3,4,5,6,7,8\}
\end{equation}

对应 golden output 为：

\begin{equation}
    y=\{1,4,9,16,25,36,49,64\}
\end{equation}

选择对角矩阵是为了让每个输出元素都具有简单且不同的 golden value，便于 UART 日志和 ILA 观测时快速判断结果是否错位。固件先运行软件 GEMV 并记录 cycle delta，再运行硬件 GEMV 并记录包含 MMIO 访问的端到端 cycle delta。两个 cycle delta 会写入加速器 debug 寄存器，便于 ILA 观察。

\section{程序镜像转换与关键问题修复}

\subsection{64-bit 程序镜像转换}

E203 在 FPGA 工程中从 ITCM 取指，而 ITCM 的仿真 RAM 宽度为 64 bit。WSL SDK 生成的 \texttt{.verilog} memory image 则是 byte-per-token 格式。二者格式不匹配：如果直接读取 byte token，Vivado 会把每个 byte 当成一个 64-bit word，导致指令排列完全错位，CPU 取到的机器码不再是原始程序。

因此程序镜像转换不是附属工具步骤，而是系统启动链路的一部分。转换逻辑按 8 byte 为一组，将 byte stream 打包为 little-endian 的 64-bit word，使 \texttt{\$readmemh} 初始化后的 ITCM 内容与 E203 取指端口宽度一致。这样 CPU 才能从地址 0 开始按预期取到固件指令。

\subsection{BRAM 推断失败问题}

上板初期，CPU 长时间停留在 \texttt{0x00000000}，说明问题不在 GEMV 计算本身，而在更早的启动或取指链路。进一步检查综合网表后发现设计中没有推断出 BRAM，意味着 ITCM 和 DTCM 没有以真实片上存储的形式存在。CPU 从异常的 ITCM 中取到无效指令后触发 trap，最终回到 \texttt{mtvec=0} 并卡死在 PC 0。

根因是原始 \texttt{sirv\_sim\_ram.v} 使用 generate 循环展开出多个 \texttt{always @(posedge clk)} 块，每个字节通道一个写块。Vivado 对这种写法无法稳定推断 block RAM。这个问题说明，硬件框架设计不仅包括计算模块，还包括程序存储、复位、时钟和综合器可识别的存储器编码风格；其中任何一环失效，CPU 都无法运行到加速器访问阶段。

修复内容如下：

\begin{itemize}
    \item 将所有字节通道写操作合并到同一个 \texttt{always @(posedge clk)} 块。
    \item 使用 \texttt{[8*j +: 8]} 完成变量 part-select，避免非法动态区间写法。
    \item 补充 \texttt{genvar i;} 声明，使底部 generate 逻辑可以综合。
    \item 为 RAM 增加 \texttt{(* ram\_style = "block" *)} 属性，辅助 Vivado 推断 BRAM。
\end{itemize}

修复后综合通过，Block RAM Final Mapping Report 中观察到的 BRAM 使用如\cref{tab:bram} 所示。

\begin{table}[H]
    \centering
    \caption{BRAM 推断结果}
    \label{tab:bram}
    \begin{tabular}{lcc}
        \toprule
        \textbf{模块} & \textbf{大小} & \textbf{RAMB36E2 数量} \\
        \midrule
        ITCM & 8K$\times$64 & 16 \\
        DTCM & 16K$\times$32 & 16 \\
        合计 & - & 32 \\
        \bottomrule
    \end{tabular}
\end{table}

该问题修复后，E203 可以从 ITCM 正常取指并执行固件。

\section{上板调试与当前验证结果}

\subsection{ILA 配置}

ILA 的作用不是替代 UART 输出，而是在硬件内部提供状态证据。早期调试阶段曾临时观察 E203 IFU 和寄存器堆，用于定位 CPU 是否正常取指；在 CPU 启动链路稳定后，当前工程保留 GEMV 外设内部 ILA，重点观察加速器的控制状态、列计数、命令触发和输出结果。

\begin{table}[H]
    \centering
    \caption{当前 GEMV ILA 观测点}
    \label{tab:ila}
    \begin{tabular}{p{3.7cm}p{3.4cm}p{6.2cm}}
        \toprule
        \textbf{观测信号} & \textbf{位置} & \textbf{作用} \\
        \midrule
        \texttt{ctrl\_done} & GEMV accelerator & 判断硬件计算是否完成 \\
        \texttt{state}, \texttt{col} & GEMV accelerator & 观察 IDLE/COMPUTE/WRDONE 状态和 8 列计算过程 \\
        \texttt{cmd\_fire}, \texttt{icb\_cmd\_addr} & ICB slave interface & 观察 CPU 是否真正访问 GEMV 地址窗口 \\
        \texttt{dbg\_hw\_cycles}, \texttt{y\_mem[0]} & GEMV accelerator & 关联端到端周期和第一个输出结果 \\
        \bottomrule
    \end{tabular}
\end{table}

GEMV 外设内部 ILA 当前连接如下，\texttt{probe0} 被打包为控制状态和调试计数，\texttt{probe1} 连接第一个输出元素：

\begin{lstlisting}[language=Verilog,caption={GEMV 外设 ILA probe}]
ila_0 u_ila_gemv (
    .clk    (clk),
    .probe0 ({ctrl_done, state, col, cmd_fire,
              icb_cmd_addr[11:0], dbg_hw_cycles[12:0]}),
    .probe1 (y_mem[0])
);
\end{lstlisting}

\subsection{PL UART 接出与串口配置}

为获得完整输出结果，本项目将 E203 内部 UART0 通过 GPIOA 的 IOF 复用接到 FPGA 顶层，并绑定到 AXU3EG 板卡的 PL UART 管脚。E203 SDK 在启动阶段执行 \texttt{gpio\_iof\_config(GPIOA, IOF\_UART\_MASK)} 和 \texttt{uart\_init(SOC\_DEBUG\_UART, 115200)}，因此固件中的 \texttt{printf} 会通过 UART0 输出。

UART0 的板端连接如\cref{tab:uart} 所示。

\begin{table}[H]
    \centering
    \caption{PL UART 管脚连接}
    \label{tab:uart}
    \begin{tabular}{lll}
        \toprule
        \textbf{E203 信号} & \textbf{顶层端口} & \textbf{FPGA 管脚} \\
        \midrule
        \texttt{GPIOA[16] / UART0 RX} & \texttt{uart0\_rxd} & \texttt{AH11} \\
        \texttt{GPIOA[17] / UART0 TX} & \texttt{uart0\_txd} & \texttt{AH12} \\
        \bottomrule
    \end{tabular}
\end{table}

上板测试时使用板卡 PL UART 对应的 CP210x USB-UART 设备，Windows 中识别为 \texttt{COM13}。串口参数为 \texttt{115200 8N1}，关闭 flow control。UART 输出用于给出完整的结果列表和 cycle 数据，ILA 用于补充硬件内部状态，两者共同构成板端验证链路。

\subsection{板端串口结果}

早期 ILA 调试中，已经通过 GEMV ILA 观察到：

\begin{lstlisting}[caption={当前 ILA 观测结果}]
ctrl_done = 1
y_mem[0]  = 0x00000001
\end{lstlisting}

最终板端测试改用 PL UART 输出完整运行日志。当前测试数据的 expected output 是：

\begin{equation}
    y=\{1,4,9,16,25,36,49,64\}
\end{equation}

串口输出摘录如下：

\begin{lstlisting}[caption={PL UART 板端输出摘录}]
HummingBird SDK Build Time: May  4 2026, 16:37:03
Download Mode: ILM
CPU Frequency 15991767 Hz
INT8 GEMV Accelerator Demo
Matrix shape: 8 x 8

SW y[0] = 1, golden = 1
SW y[1] = 4, golden = 4
SW y[2] = 9, golden = 9
SW y[3] = 16, golden = 16
SW y[4] = 25, golden = 25
SW y[5] = 36, golden = 36
SW y[6] = 49, golden = 49
SW y[7] = 64, golden = 64
SW GEMV PASS
SW Cycle delta: 1630

HW y[0] = 1, golden = 1
HW y[1] = 4, golden = 4
HW y[2] = 9, golden = 9
HW y[3] = 16, golden = 16
HW y[4] = 25, golden = 25
HW y[5] = 36, golden = 36
HW y[6] = 49, golden = 49
HW y[7] = 64, golden = 64
HW GEMV PASS
HW Cycle delta: 243
\end{lstlisting}

该结果说明：

\begin{itemize}
    \item E203 固件已经能够运行到访问 GEMV 外设的位置。
    \item CPU 可以通过 MMIO 访问 \texttt{0x1004\_1000} 地址段。
    \item CPU 已经成功写入数据并启动加速器。
    \item 加速器已经完成计算，并产生了完整正确的 8 个输出元素。
    \item 软件 GEMV 板端执行周期为 1630 cycles。
    \item 硬件加速器端到端调用周期为 243 cycles。
\end{itemize}

这表明从 CPU 到外设的控制链路、数据链路和串口观测链路均已打通。ILA 结果用于证明硬件内部状态确实发生变化，PL UART 结果用于给出完整正确性和性能数据。

\subsection{Reset 与时钟设计}

旧设计中板上 reset 会复位 MMCM，导致 ILA/debug hub 依赖的 \texttt{clk\_16M} 停止，Vivado Hardware Manager 可能报无法连接 debug core。当前已修改 reset 方案：

\begin{itemize}
    \item MMCM \texttt{resetn} 固定为 \texttt{1'b1}，不再由板上 reset 拉停。
    \item CPU reset 改为由 \texttt{mmcm\_locked} 和上电延迟释放逻辑控制。
    \item 上电延迟计数宽度为 20 bit，在 16 MHz 下约为 65 ms。
\end{itemize}

这个修改体现了板端调试时的一个设计原则：调试基础设施依赖的时钟不能被普通业务 reset 随意关闭。当前设计将 MMCM 保持运行，只对 CPU 复位释放做延迟控制，从而避免 reset 操作破坏 ILA/debug hub 连接。

\section{性能分析}

\subsection{核心计算与端到端性能}

软件基线、硬件端到端调用和 RTL 核心计算周期对比如\cref{tab:perf-core} 所示。其中，软件和硬件端到端数据来自 PL UART 板端输出；RTL 核心计算延迟来自 \texttt{gemv\_accel.v} 的 FSM 结构推导。

\begin{table}[H]
    \centering
    \caption{软件、硬件端到端与 RTL 核心周期对比}
    \label{tab:perf-core}
    \begin{tabular}{p{4.0cm}p{5.7cm}p{3.2cm}}
        \toprule
        \textbf{实现方式} & \textbf{计算内容} & \textbf{周期数} \\
        \midrule
        E203 软件 kernel & 8$\times$8 INT8 GEMV 软件 & 1630 cycles \\
        GEMV 硬件端到端 & 写入 $W/x$、启动、轮询 done、读取 $y$ & 243 cycles \\
        GEMV RTL 核心 & 8 路并行 MAC 计算与结果写回 & 9 cycles \\
        \bottomrule
    \end{tabular}
\end{table}

硬件端到端加速比为：

\begin{equation}
    S_{\text{end-to-end}} = \frac{1630}{243} \approx 6.71
\end{equation}

硬件核心只需要 9 cycles 的原因是它将 8 个输出元素的 MAC 并行展开，每个周期处理一个输入列。软件实现虽然也只计算 64 次乘加，但每次乘加、循环控制、访存和寄存器调度都由 E203 顺序执行，因此 cycle 数明显更高。

需要注意，9 cycles 是 RTL 状态机推导出的核心计算延迟，并非板上串口直接测得的端到端周期。板上实测的 243 cycles 更能反映 CPU 实际调用加速器时的系统级开销。

\subsection{端到端开销}

实际系统中，CPU 使用硬件加速器时不仅要等待核心计算，还需要完成以下操作：

\begin{itemize}
    \item 写入 16 个 32-bit word 的矩阵数据。
    \item 写入 2 个 32-bit word 的输入向量。
    \item 写 \texttt{CTRL} 启动外设。
    \item 轮询 \texttt{CTRL.done}。
    \item 读取 8 个 32-bit 输出结果。
\end{itemize}

因此硬件辅助端到端周期一定大于 9 cycles。最终报告中将三类指标分开呈现：

\begin{enumerate}
    \item 软件 GEMV kernel cycle：1630 cycles。
    \item GEMV RTL 核心计算 cycle：9 cycles，由 FSM 推导。
    \item CPU + MMIO + GEMV 外设的硬件辅助端到端 cycle：243 cycles。
\end{enumerate}

从端到端结果看，硬件加速器虽然核心计算只需 9 cycles，但一次调用中仍包含 16 次权重写入、2 次输入向量写入、1 次 start 写入、done 轮询和 8 次输出读取。对于 8$\times$8 这样较小的矩阵，MMIO 访问开销占比较高，因此端到端加速比低于核心计算理论加速比。若矩阵规模扩大或采用批量数据搬运方式，硬件核心并行度带来的收益会更明显。

\section{当前限制与后续工作}

目前项目已经完成主体实现和完整上板验证，但仍存在以下可改进点：

\begin{enumerate}
    \item 若需要更强的板上波形证据，可使用当前 GEMV ILA 抓取 \texttt{state}、\texttt{col} 和 \texttt{ctrl\_done}，展示 8 个 COMPUTE 周期和 1 个 WRDONE 周期。
    \item 当前 GEMV 已将 8-bit operand sign-extend 到 16-bit 后再相乘，后续可继续约束乘法器资源映射，降低综合资源和综合时间。
    \item 当前工程只保留 GEMV 相关 ILA。若最终演示完全依赖 UART，可以进一步生成无 ILA 的发布 bitstream，减少 debug 资源占用。
    \item 当前测试规模为 8$\times$8，适合验证链路。后续可扩展为更大矩阵或 tile-based GEMV，以提高硬件计算时间在端到端开销中的占比。
\end{enumerate}

\section{结论}

本项目围绕 INT8 GEMV 算子，完成了 E203 软件仿真、FPGA 工程迁移、硬件加速器设计、SoC 外设集成和完整上板验证。软件阶段建立了 8$\times$8 INT8 GEMV 的正确性和性能基线，早期 iVerilog 仿真测得软件 kernel 为 1638 cycles，最终板端 UART 测得软件 GEMV 为 1630 cycles。硬件阶段实现了一个挂接在 E203 peripheral ICB bus 上的 memory-mapped GEMV 加速器，采用 8 路并行 signed MAC，可在 RTL 核心中以 9 cycles 完成 8$\times$8 GEMV 的核心计算。

调试过程中，项目定位并修复了 \texttt{sirv\_sim\_ram.v} 无法被 Vivado 推断为 BRAM 的关键问题，使 E203 能够在 FPGA 上从 ITCM 正常取指运行。早期 ILA 证明 CPU 可以通过 MMIO 启动 GEMV 外设，最终 PL UART 输出进一步证明 8 个硬件输出元素全部与 golden values 一致，\texttt{SW GEMV PASS} 和 \texttt{HW GEMV PASS} 均成立。板端实测硬件端到端调用周期为 243 cycles，相比软件 1630 cycles 达到约 $6.71\times$ 端到端加速。

\end{document}
